\documentclass[12pt,letterpaper]{article}
\usepackage[T1]{fontenc}
\usepackage{amsmath,amsthm}
\usepackage{newtxtext,newtxmath}
\usepackage{array,booktabs}
\usepackage[letterpaper,margin=1in]{geometry}
\usepackage{url}
\usepackage[hidelinks]{hyperref}

\newtheorem{theorem}{Theorem}[section]
\newtheorem{corollary}[theorem]{Corollary}
\newtheorem{conjecture}[theorem]{Conjecture}
\newtheorem{lemma}[theorem]{Lemma}
\newtheorem{proposition}[theorem]{Proposition}
\theoremstyle{definition}
\newtheorem{definition}[theorem]{Definition}
\newtheorem{example}[theorem]{Example}
\newtheorem{problem}[theorem]{Problem}
\theoremstyle{remark}
\newtheorem{remark}[theorem]{Remark}
\numberwithin{equation}{section}

\newcommand{\F}{\mathbb{F}}
\newcommand{\PP}{\mathbb{P}}
\newcommand{\supp}{\operatorname{supp}}
\newcommand{\Aut}{\operatorname{Aut}}
\newcommand{\Stab}{\operatorname{Stab}}
\newcommand{\ev}{\operatorname{ev}}
\newcommand{\wt}{\operatorname{wt}}
\newcommand{\RS}{\operatorname{RS}}
\newcommand{\PAut}{\operatorname{PAut}}
\newcommand{\AGL}{\operatorname{AGL}}
\newcommand{\GL}{\operatorname{GL}}
\newcommand{\PSL}{\operatorname{PSL}}
\newcommand{\PGL}{\operatorname{PGL}}
\newcommand{\ord}{\operatorname{ord}}
\newcommand{\Gsys}{\mathsf{G}}
\newcommand{\Hsys}{\mathsf{H}}
\newcommand{\cI}{\mathcal{I}}
\newcommand{\cJ}{\mathcal{J}}

\begin{document}

\title{Conjectural permutation decoding of some AG codes\thanks{This is a
corrected and expanded version of the note~\cite{Jo05}, published in the
\emph{ACM SIGSAM Bulletin} in 2005. This version, Draft~2, is dated
November 12, 2026.}\\[4pt]
\large Corrected and expanded version. Draft 2}
\author{David Joyner\thanks{Original author. \texttt{wdjoyner@gmail.com}}
\and ChatGPT (OpenAI)\thanks{AI system: first mathematical revision,
proof checking, and exposition.  Listed as an author at the request
of the first author.}
\and Claude (Anthropic)\thanks{AI system: refereeing, independent
verification, exhaustive computations, literature review, and preparation
of this version.  Listed as an author at the request of the first author.
See the Acknowledgments for a description of each contribution.}}
\date{Original draft: 2005. Corrected and expanded draft: 2026.}
\maketitle

\begin{abstract}
This paper was revised with the help of two AI systems, ChatGPT (OpenAI)
and Claude (Anthropic), which are also listed as co-authors. The
Acknowledgments give further details.
The paper revisits the permutation-decoding conjectures of
the first author (2005) for algebraic-geometric codes associated with the
hyperelliptic curves $y^2=x^p-x$.  The one-point codes of length $p$ over
$\F_p$ turn out to be full-length Reed--Solomon codes, whose permutation
group is the affine group $\AGL(1,p)$.  This group always gives a
$2$-error permutation-decoding set.  Exhaustive computation shows that for
$p\le 23$ it is a full permutation-decoding set for a suitable choice of
information set.  However, at $p=17$ some information sets fail even
though the standard counting bound is satisfied, and at $p=37$ the
counting bound rules out full permutation decoding for every information
set.  For the codes of length $2p(p-1)$ over $\F_{p^2}$
($p\equiv3\pmod4$), the stabilizer of the point at infinity was proposed
as a permutation-decoding set.  It acts regularly and is too small: we
give an infinite family of half-rate codes for which it is not even a
$(\lfloor\log_2 n\rfloor+1)$-error decoding set.  We also correct several
statements about automorphism groups, separation of points, and
permutation modules.  In particular, we exhibit a six-dimensional
Riemann--Roch module that is a submodule of a permutation module but is
not itself a permutation module.  A background section, a literature
review, and a glossary are included.
\end{abstract}

\tableofcontents

%======================================================================
\section{Introduction}\label{sec:intro}

This section states the problem and recalls the two conjectures of the
note~\cite{Jo05}. It then gives a small example, summarizes what this
paper proves about the conjectures, and describes how the paper is
organized.

Let $X$ be a smooth projective curve over a finite field $F$, let $D$ be
a divisor on $X$, and let $E=P_1+\cdots+P_n$ be a sum of distinct
$F$-rational points not in the support of $D$.  The algebraic-geometric
(AG) code $C_X(D,E)$ is obtained by evaluating the functions in the
Riemann--Roch space $L(D)$ at $P_1,\dots,P_n$.  If a group $G$ of
automorphisms of $X$ preserves both $D$ and $E$, then $G$ permutes the
coordinates of $C_X(D,E)$ and preserves the code.  Permutation decoding,
introduced by Prange and MacWilliams \cite{Pr62,Mac64}, is a natural way
to exploit such symmetry.  A set $S$ of code automorphisms is used to move
every error pattern of bounded weight away from a fixed information set,
where it can then be removed by re-encoding.

The note \cite{Jo05}, of which this paper is a corrected and expanded
version, proposed that the curves
\begin{equation}\label{eq:curve}
X:\qquad y^2=x^p-x,
\end{equation}
whose automorphism groups are unusually large, give AG codes with very
efficient permutation decoding.  It formulated two conjectures:
\begin{itemize}
\item[(C1)] for one-point codes of length $p$ over $\F_p$, permutation
decoding always applies, at a cost of at most $|\PAut(C)|=O(p^2)$ codeword
operations;
\item[(C2)] for one-point codes of length $2p(p-1)$ over $\F_{p^2}$, with
$p\equiv 3\pmod 4$, the image of the stabilizer of the point at infinity
can serve as a full permutation-decoding set.
\end{itemize}
According to \cite{LLMS26}, these conjectures were essentially the only
prior work on permutation decoding of AG codes of positive genus.  This
paper settles both conjectures in the negative, as universal statements,
and replaces them with precise positive results, explicit obstructions,
and open problems.

\begin{example}\label{ex:intro-p7}
Let $p=7$. The curve $y^2=x^7-x$ has genus $3$. Over $\F_7$ its affine
rational points are the seven points $(a,0)$ with $a\in\F_7$. Evaluating
the functions $1$, $x$ and $x^2$ at these points gives a code of length
$7$ and dimension $3$ over $\F_7$. This is the $[7,3,5]$ code studied
in~\cite{Jo05}. It corrects two errors, and its permutation group has
order $42$. The whole group is a full permutation-decoding set for this
code (Example~\ref{ex:7-3-5}).
\end{example}

\paragraph{Summary of results.}
\begin{enumerate}
\item The codes in (C1) are full-length Reed--Solomon codes, with
permutation group $\AGL(1,p)$ (Theorem~\ref{thm:RS}).
\item $\AGL(1,p)$ is a $2$-PD-set for every information set
(Theorem~\ref{thm:two-pd}).
\item For $p\le23$, $\AGL(1,p)$ is a full PD-set relative to some
information set, in every dimension.  For $p\le13$ every information set
works, but from $p=17$ on some do not, even though the counting
bound holds (Proposition~\ref{prop:small-p}, Example~\ref{ex:p17}).
\item At $p=37$ no information set works (Theorem~\ref{thm:counter-RS}).
Hence (C1) fails.
\item Over $\F_{p^2}$, a regular group of coordinate permutations cannot
correct more than about $\log n/\log(n/(n-k))$ errors by permutation
decoding (Theorem~\ref{thm:regular}).  This gives an infinite family of
counterexamples to (C2) (Corollary~\ref{cor:family}).
\item The representation of the automorphism group on $L(D)$ need not be
a permutation representation, even under the strong degree hypotheses
used in \cite{Jo05} (Example~\ref{ex:not-permutation}).
\end{enumerate}

\paragraph{Relation to the original note.}
We keep the curve-theoretic setup, the GUAVA example, and the research
motivation of \cite{Jo05}.  Several statements there needed correction.
They are flagged where they occur, notably in
Remark~\ref{rem:original-notation} (point labels),
Section~\ref{subsec:fp} (stabilizer orders and the action of the
hyperelliptic involution), Remark~\ref{rem:separation} (ampleness versus
separation of points), and Section~\ref{subsec:not-perm} (permutation
modules).

\paragraph{Organization.}
Section~\ref{sec:lit} reviews the literature before and after 2007.
Section~\ref{sec:background} recalls the background needed later and ends
with a table of notation.
Sections~\ref{sec:curve}--\ref{sec:fp2-codes} contain the mathematics,
Section~\ref{sec:complexity} discusses complexity, and
Section~\ref{sec:problems} summarizes the results and lists open problems.
The appendices contain
the historical GUAVA session and a short verification script.  A glossary
of technical terms precedes the Acknowledgments and the references.

%======================================================================
\section{Literature review}\label{sec:lit}

This section places the paper in context. It is split at 2007 only for
convenience. The first subsection covers work up to the time of the
note~\cite{Jo05} and the two years after it. The second covers later
work.

\subsection{Work up to 2007}

The earlier work falls into four strands, treated in turn: permutation
decoding, AG codes and their automorphisms, representations on
Riemann--Roch spaces, and the decoding of AG codes.

\paragraph{Permutation decoding.}
Prange introduced the use of information sets to decode cyclic codes
\cite{Pr62}, and MacWilliams developed permutation decoding for general
systematic codes \cite{Mac64}.  A textbook account is in MacWilliams and
Sloane \cite[Ch.~16]{MS77}.  Early applications treated specific cyclic
codes \cite{SFT70,SF72}.  Gordon gave a method for finding minimal
decoding sets for the binary Golay codes \cite{Gor82}, together with the
covering-type lower bound that underlies Theorem~\ref{thm:counting}
below.  Wolfmann found a set of fourteen permutations for the
$[24,12,8]$ Golay code \cite{Wo83}.  Chabanne extended permutation
decoding to abelian codes using Gr\"obner bases \cite{Ch92}, and Biglieri
considered group codes \cite{Bi95}.  Huffman's survey chapter \cite{Hu98}
and the textbook \cite{HP03} summarize the theory.

The notion of \emph{partial} permutation decoding, using $s$-PD-sets with
$s<t$, was introduced by Key, McDonough and Mavron \cite{KMM05} and
also by Kroll and Vincenti \cite{KV05}.  Permutation decoding and its
partial version were applied to codes from designs, finite geometries and
graphs \cite{KMR04,KMM06,KS07}.  These works emphasize two points that are
central here: the choice of information set matters, and a large
permutation group does not by itself guarantee full permutation decoding.

\paragraph{AG codes and their automorphisms.}
AG codes were introduced by Goppa \cite{Go81}.  The Tsfasman--Vl\u{a}du\c{t}--Zink
theorem \cite{TVZ82} showed that they can beat the Gilbert--Varshamov
bound.  Standard references are \cite{St93,TV91,HvLP98}.  The relation between
curve automorphisms and code automorphisms was studied by Stichtenoth
\cite{St90} and by Wesemeyer \cite{We98} for plane curves.  Joyner and Ksir
\cite{JK06} showed that, in many cases, $\Aut(X)$ and the permutation
group of a corresponding AG code coincide.  For genus zero, the
automorphism groups of (generalized, doubly extended) Reed--Solomon codes
were determined by D\"ur \cite{Du87}.  The automorphism groups of
generalized Reed--Muller codes over prime fields, which include the
full-length Reed--Solomon codes over $\F_p$, were determined by Kn\"orr and
Willems \cite{KW90} and by Berger and Charpin \cite{BC93}.  Automorphism
groups of hyperelliptic curves were classified by Brandt and Stichtenoth
\cite{BrS86}, and computed algorithmically by G\"ob \cite{Go03}.

\paragraph{Representations on Riemann--Roch spaces.}
The action of $\Aut(X)$ on $L(D)$ was studied by Joyner and Traves
\cite{JT02} and by Joyner and Ksir \cite{JK03}, building on modular
representation theory such as \cite{Sc90}.  The conjectures on permutation
decoding of codes from $y^2=x^p-x$ first appeared in \cite{JT02} and were
formulated in \cite{Jo05}.  GUAVA 2.0 implemented permutation decoding
\cite{GUAVA}.  Shaska and Wang \cite{ShW07} studied automorphism groups of
AG codes from $C_{a,b}$ curves, motivated in part by the decoding
suggestion of Joyner and Ksir.

\paragraph{Decoding AG codes.}
Algebraic decoding of AG codes goes back to Justesen et al.\
\cite{JLJHH89}, Skorobogatov and Vl\u{a}du\c{t} \cite{SV90}, and Feng and
Rao \cite{FR93}.  List decoding was developed by Sudan \cite{Su97},
Guruswami and Sudan \cite{GS99}, and Shokrollahi and Wasserman
\cite{SW99}.  The survey \cite{HvLP98} describes the state of the art around
1998.  These algorithms decode up to (or beyond) half the designed
distance in polynomial time.  The interest of permutation decoding lies
in its simplicity: it uses only syndrome computations and coordinate
permutations.

\subsection{Work after 2007}

Later work has developed partial permutation decoding and clarified the
permutation groups of Reed--Solomon codes. It has also returned to
permutation decoding of AG codes of positive genus, which bears most
directly on this paper.

\paragraph{Partial permutation decoding.}
Partial PD-sets were constructed for codes from lattice graphs
\cite{KS08}, for first-order Reed--Muller codes \cite{Sen09}, and for
codes from affine geometry designs \cite{KMM08}.  Key, McDonough and
Mavron further studied Reed--Muller codes \cite{KMM10,KMM17}.  Lim,
Fossorier and Kav\v{c}i\'c used code automorphisms for permutation
decoding of binary images of Reed--Solomon codes \cite{LFK10}.  Bernal and
Sim\'on developed a systematic theory of information sets for abelian
codes in terms of defining sets \cite{BS11,BS14}, and used it for partial
permutation decoding of abelian and affine-invariant codes
\cite{BS13,BS23}.

\paragraph{Permutation groups of Reed--Solomon codes.}
Camps-Moreno, Lau, L\'opez and Santos gave an elementary proof \cite{CLLS25}
that the permutation group of a Reed--Solomon code on an arbitrary
evaluation set consists of the affine maps preserving that set, away from
exceptional dimensions.  This recovers the case of the full field used in
Theorem~\ref{thm:RS}.  Joyner, Ksir and Traves \cite{JKT08} revisited
automorphism groups of AG codes on $\PP^1$ using algebraic-geometric
methods.

\paragraph{Permutation decoding of AG codes of positive genus.}
Recent work has returned to the question raised in \cite{Jo05}.
Lichtenwalner, L\'opez, Matthews and Seneviratne \cite{LLMS26} use
automorphisms of Hermitian and norm-trace curves to construct PD-sets that
correct specific \emph{burst} errors for one-point codes.  They note that
little progress had been made for positive genus beyond the conjectures
of \cite{Jo05}.  Castellanos, Tizziotti and Olaya-Le\'on \cite{CTO26}
treat curves defined by separated polynomials $f(y)=g(x)$ with $f$
additive, including Hermitian and generalized Hermitian curves.  They
obtain PD-sets for bursts on points sharing a coordinate, and partial
$2$-PD-sets formed from two families of automorphisms.

Note that \eqref{eq:curve} can be written $x^p-x=y^2$, which has this
separated form with the additive polynomial in $x$.  Its Weierstrass
semigroup at $P_\infty$ is generated by $2$ and $p$, the degrees of the
two sides.  The $2$-PD result of Theorem~\ref{thm:two-pd} below is analogous to
their two-family construction: the translations $x\mapsto x+b$ are
exactly the automorphisms arising from the roots of the additive
polynomial $x^p-x$.  Whether the burst-error PD-sets of \cite{CTO26,LLMS26} apply to
the codes of Section~\ref{sec:fp2-codes} is raised in
Problem~\ref{prob:burst}.

%======================================================================
\section{Background}\label{sec:background}

This section collects standard facts, with references, so that the rest
of the paper can be read with modest prerequisites.  Proofs are omitted.
Each subsection ends with an example, and all of these examples use the
prime $p=7$. The section closes with a table of notation.

\subsection{Linear codes}\label{subsec:bg-codes}
This subsection fixes the vocabulary of linear codes: weight, minimum
distance, information sets, systematic matrices and permutation
automorphisms. Permutation decoding is defined in these terms in
Section~\ref{sec:PD}.

A \emph{linear code} $C$ of length $n$ over $F=\F_q$ is a subspace of
$F^n$.  The \emph{Hamming weight} $\wt(v)$ of $v\in F^n$ is the number of
nonzero coordinates of $v$.  The \emph{minimum distance} $d$ of $C\neq0$
is the least weight of a nonzero codeword.  We say $C$ is an $[n,k,d]$
code if $\dim C=k$.  Such a code corrects any $t=\lfloor(d-1)/2\rfloor$
errors.  The Singleton bound $d\le n-k+1$ holds, and codes attaining it
are called \emph{maximum distance separable} (MDS) \cite{HP03,MS77}.

An \emph{information set} is a set $\cI$ of $k$ coordinates such that
projection $C\to F^{\cI}$ is an isomorphism.  The complementary set
$\cJ$ of $r=n-k$ coordinates is the set of \emph{check positions}.  For
an MDS code, every $k$-subset is an information set \cite[Ch.~11]{MS77}.
After reordering coordinates so that $\cI$ comes first, $C$ has a
\emph{systematic} generator matrix $\Gsys=(I_k\mid A)$ and parity-check
matrix $\Hsys=(-A^{\mathsf T}\mid I_r)$.  The \emph{syndrome} of
$v\in F^n$ is $\Hsys v^{\mathsf T}$.

The symmetric group $S_n$ acts on $F^n$ by
$(\sigma v)_i=v_{\sigma^{-1}(i)}$.  The \emph{permutation automorphism
group} of $C$ is
$\PAut(C)=\{\sigma\in S_n:\sigma(C)=C\}$.
(Monomial automorphisms, which also rescale coordinates, will not be
needed.)

\begin{example}\label{ex:bg-codes}
The matrix $\Gsys$ of Example~\ref{ex:7-3-5} generates a $[7,3,5]$ code
$C$ over $\F_7$. It has the form $(I_3\mid A)$, so the first three
coordinates form an information set $\cI$ and the last four are the check
positions $\cJ$. Here $r=4$ and $t=2$. Since $5=7-3+1$, the code is MDS,
and every $3$-subset of the coordinates is an information set. The
parity-check matrix is
\[
\Hsys=\begin{pmatrix}
1&1&4&1&0&0&0\\
3&4&6&0&1&0&0\\
1&4&1&0&0&1&0\\
3&6&4&0&0&0&1
\end{pmatrix}.
\]
The first row $c=(1,0,0,6,4,6,4)$ of $\Gsys$ has weight $5$ and syndrome
$\Hsys c^{\mathsf T}=0$.
\end{example}

\subsection{Reed--Solomon codes and the affine group}\label{subsec:bg-rs}
The codes of length $p$ studied in Section~\ref{sec:RS} turn out to be
Reed--Solomon codes. This subsection recalls their definition and the
affine group, which permutes their coordinates.

For $1\le k\le q$, the full-length Reed--Solomon code
\[
\RS(\F_q,k)=\{(f(a))_{a\in\F_q}:\ f\in\F_q[x],\ \deg f<k\}
\]
is a $[q,k,q-k+1]$ MDS code.  The \emph{affine group}
$\AGL(1,q)=\{x\mapsto ux+b:\ u\in\F_q^\times,\ b\in\F_q\}$, of order
$q(q-1)$, acts on the coordinates by substitution and preserves every
$\RS(\F_q,k)$.  It is \emph{sharply $2$-transitive}: for distinct
$a_1,a_2$ and distinct $b_1,b_2$ there is exactly one affine map with
$a_i\mapsto b_i$.

\begin{example}\label{ex:rs7}
Let $q=7$ and $k=3$. The polynomial $x^2$ gives the codeword
$(0,1,4,2,2,4,1)$ of $\RS(\F_7,3)$, listed in the order $a=0,1,\dots,6$.
The polynomial $(x-1)(x-2)$ gives $(2,0,0,2,6,5,6)$, which has weight
$5=7-3+1$. The group $\AGL(1,7)$ has order $42$. The unique affine map
with $0\mapsto3$ and $1\mapsto5$ is $x\mapsto2x+3$.
\end{example}

\subsection{Curves, divisors, and Riemann--Roch spaces}\label{subsec:bg-curves}
AG codes are built from spaces of functions on a curve with prescribed
poles. This subsection recalls divisors, Riemann--Roch spaces and the
evaluation codes built from them.

We follow Stichtenoth \cite{St93} and Hartshorne \cite{Ha77}.  Let $X/F$ be a
smooth, projective, geometrically irreducible curve of genus $g$.  A
\emph{divisor} is a finite formal sum $D=\sum n_PP$ of points, with
degree $\deg D=\sum n_P\deg P$.  For a nonzero function $f$,
$\operatorname{div}(f)=\sum\ord_P(f)P$ has degree $0$.  The
\emph{Riemann--Roch space} is
\[
L(D)=\{f\in F(X)^\times:\ \operatorname{div}(f)+D\ge0\}\cup\{0\},
\qquad \ell(D)=\dim_F L(D).
\]
The Riemann--Roch theorem gives $\ell(D)\ge\deg D+1-g$, with equality
when $\deg D>2g-2$ \cite[Thm.~I.5.17]{St93}.  If
$E=P_1+\dots+P_n$ consists of rational points outside $\supp D$, the
evaluation map $\ev_E:L(D)\to F^n$, $f\mapsto(f(P_i))_i$, has image the
\emph{AG code} $C_X(D,E)$ \cite{Go81,TVZ82}.  If $\deg D<n$ then $\ev_E$ is
injective, and $C_X(D,E)$ is an $[n,\ell(D),\ge n-\deg D]$ code
\cite[Thm.~II.2.2]{St93}.  When $D=mQ$ for a single point $Q$, the code is
called a \emph{one-point code}.

\begin{example}\label{ex:bg-p1}
Let $X=\PP^1$ over $\F_7$, of genus $0$, with coordinate $x$, and let
$P_\infty$ be the pole of $x$. For $D=2P_\infty$ we get
$L(D)=\langle1,x,x^2\rangle$ and $\ell(D)=3=\deg D+1-g$. Let $E$ be the
sum of the seven points $x=a$ with $a\in\F_7$. Then
$C_X(D,E)=\RS(\F_7,3)$. Since $\deg D=2<7$, the evaluation map is
injective and the minimum distance is at least $7-2=5$. This matches
Example~\ref{ex:rs7}.
\end{example}

\subsection{Hyperelliptic curves and point counts}\label{subsec:bg-hyper}
The curves of this paper are hyperelliptic. This subsection recalls the
hyperelliptic involution, the Weierstrass points and the Hasse--Weil
bound on the number of rational points.

A curve of genus $g\ge2$ is \emph{hyperelliptic} if it admits a degree-two
map $x:X\to\PP^1$.  The associated \emph{hyperelliptic involution}
$\iota$ is central in $\Aut(X)$.  Its fixed points are the
\emph{Weierstrass points}, which are the ramification (branch) points of
$x$.  Every automorphism of $X$ therefore induces an automorphism of
$\PP^1$ that permutes the branch locus \cite{BrS86,Go03}.  Over $F=\F_q$, the
\emph{Hasse--Weil bound} says
\[
|\,|X(\F_q)|-(q+1)\,|\le 2g\sqrt q.
\]
Curves attaining the upper (respectively lower) bound are called
\emph{maximal} (respectively \emph{minimal}) \cite[Thm.~V.2.3]{St93}.

\begin{example}\label{ex:bg-hyper}
Let $X$ be the curve $y^2=x^7-x$ over $\F_7$, which has genus $3$
(Section~\ref{sec:curve}). The map $x$ has degree two, and $\iota$ is
$(x,y)\mapsto(x,-y)$. The Weierstrass points are the seven points $(a,0)$
with $a\in\F_7$ and the point at infinity. That makes $8=2g+2$ points.
These are also all the $\F_7$-rational points, and $|8-(7+1)|=0$ lies
within the Hasse--Weil bound $2\cdot3\sqrt7\approx15.9$. Over $\F_{49}$
the curve has $92$ rational points (Proposition~\ref{prop:orbits}). Since
$92=49+1+2\cdot3\cdot7$, the curve is maximal over $\F_{49}$.
\end{example}

\subsection{Group actions and modules}\label{subsec:bg-groups}
Permutation decoding rests on how a group moves coordinates, and
Section~\ref{sec:modules} compares two kinds of modules. This subsection
recalls orbits, regular actions, permutation modules and the convention
for the action of a group on functions.

A finite group $G$ acting on a finite set $\Omega$ satisfies
$|G\cdot\omega|=|G|/|G_\omega|$ (orbit--stabilizer).  The action is
\emph{regular} if it is transitive with trivial stabilizers, in which
case $|G|=|\Omega|$.  The \emph{permutation module} $F^\Omega$ has basis
$\Omega$ permuted by $G$.  A submodule of a permutation module need not be
a permutation module.  A finite-dimensional $F[G]$-module $V$ has a
composition series, and the direct sum of its composition factors is its
\emph{semisimplification}.

If $G\le\Aut_F(X)$ preserves $D$, then $G$ acts on $L(D)$.  Throughout
this paper we use the convention
\begin{equation}\label{eq:convention}
(g\cdot f)(P)=f(g^{-1}P),\qquad\text{i.e. }g\cdot f=f\circ g^{-1}.
\end{equation}

\begin{example}\label{ex:bg-groups}
Let $G=\AGL(1,7)$ act on $\Omega=\F_7$. The action is transitive. The
stabilizer of $0$ consists of the six maps $x\mapsto ux$, and $42/6=7$,
as orbit--stabilizer requires. The action is not regular. The subgroup of
the seven translations $x\mapsto x+b$ does act regularly.

In the permutation module $F^\Omega$ with $F=\F_7$, the vectors with
coordinate sum $0$ form a submodule $W$ of dimension $6$. The translation
$x\mapsto x+1$ has order $7$ and acts nontrivially on $W$. A permutation
of six basis vectors whose order divides $7$ is the identity. So $W$ is
not a permutation module for $G$. Example~\ref{ex:not-permutation} uses
the same argument.

For the convention~\eqref{eq:convention}, let $g$ be the translation
$x\mapsto x+1$ and let $f$ be the coordinate function $x$. Then
$g\cdot f=x-1$.
\end{example}

\subsection{Projective linear groups}\label{subsec:bg-pgl}
The automorphisms of the curve are governed by the groups $\PGL_2(p)$ and
$\PSL_2(p)$ (Section~\ref{sec:curve}). This subsection recalls their
orders and their action on the projective line.

$\PGL_2(p)=\GL_2(\F_p)/\F_p^\times$ has order $p(p^2-1)$ and acts sharply
$3$-transitively on $\PP^1(\F_p)$.  Its subgroup $\PSL_2(p)$ of index two
consists of the classes of matrices with square determinant.  The
stabilizer of $\infty$ in $\PGL_2(p)$ is the \emph{Borel subgroup}
$B\cong\AGL(1,p)$ of upper-triangular classes.

\begin{example}\label{ex:bg-pgl}
For $p=7$, $|\PGL_2(7)|=7\cdot48=336$ and $|\PSL_2(7)|=168$. The
projective line $\PP^1(\F_7)$ has $8$ points, and the Borel subgroup $B$
has order $336/8=42=|\AGL(1,7)|$. The squares in $\F_7^\times$ are
$1,2,4$. So the map $x\mapsto2x$, the class of a diagonal matrix of
determinant $2$, lies in $\PSL_2(7)$. The map $x\mapsto3x$ has
determinant $3$ and lies in $\PGL_2(7)$ but not in $\PSL_2(7)$.
\end{example}

\subsection{Table of notation}\label{subsec:notation}
This subsection collects the main symbols of the paper, each with the
place where it is introduced. The glossary that precedes the
Acknowledgments explains the technical terms.

\begin{center}\small
\begin{tabular}{@{}l>{\raggedright\arraybackslash}p{0.52\textwidth}l@{}}
\toprule
Symbol & Meaning & Introduced in\\
\midrule
$p$, $g$ & a prime $p\ge5$ and the genus $g=(p-1)/2$ & Section~\ref{sec:curve}\\
$X$ & the smooth projective model of $y^2=x^p-x$ & \eqref{eq:curve}\\
$P_\infty$, $P_a$ & the point at infinity and the point $(a,0)$, $a\in\F_p$ & Section~\ref{sec:curve}\\
$\iota$, $\delta_a$, $\tau_b$, $\omega$ & explicit automorphisms of $X$ & \eqref{eq:generators}\\
$G_p$, $H_p$ & $\Aut_{\F_p}(X)$ and the stabilizer of $P_\infty$ in it & Section~\ref{subsec:fp}\\
$G_{p^2}$, $H$ & the automorphism group over $\F_{p^2}$ and the stabilizer of $P_\infty$ in it & Section~\ref{subsec:fp2}\\
$O_1$, $O_2$ & the orbits of $G_{p^2}$ on $X(\F_{p^2})$ when $p\equiv3\pmod4$ & Proposition~\ref{prop:orbits}\\
$E_2$ & the sum of the points of $O_2$ & Section~\ref{sec:fp2-codes}\\
$L(D)$, $\ell(D)$ & the Riemann--Roch space of a divisor $D$ and its dimension & Section~\ref{subsec:bg-curves}\\
$\ev_E$, $C_X(D,E)$ & evaluation at the points of $E$ and the AG code, its image & Section~\ref{subsec:bg-curves}\\
$g\cdot f$ & the function $f\circ g^{-1}$ & \eqref{eq:convention}\\
$\pi_E$, $\rho_C$ & the action of a group on coordinates and its restriction to the code & Section~\ref{sec:modules}\\
$[n,k,d]$, $t$, $r$ & length, dimension, minimum distance, $t=\lfloor(d-1)/2\rfloor$ and $r=n-k$ & Section~\ref{subsec:bg-codes}\\
$\wt$ & Hamming weight & Section~\ref{subsec:bg-codes}\\
$\cI$, $\cJ$ & an information set and the complementary check positions & Section~\ref{subsec:bg-codes}\\
$\Gsys$, $\Hsys$ & systematic generator and parity-check matrices & Section~\ref{subsec:bg-codes}\\
$\PAut(C)$ & the permutation automorphism group of $C$ & Section~\ref{subsec:bg-codes}\\
$\RS(\F_q,k)$, $\AGL(1,q)$ & the full-length Reed--Solomon code and the affine group & Section~\ref{subsec:bg-rs}\\
$\PGL_2(p)$, $\PSL_2(p)$, $B$ & the projective linear groups and the Borel subgroup & Section~\ref{subsec:bg-pgl}\\
$s$-PD-set & a set of automorphisms that moves any $s$ errors into $\cJ$ & Definition~\ref{def:pd-set}\\
(C1), (C2) & the two conjectures of~\cite{Jo05} & Section~\ref{sec:intro}\\
\bottomrule
\end{tabular}
\end{center}

%======================================================================
\section{The curve and its automorphisms}\label{sec:curve}

This section describes the curve $X$, its rational points and its
automorphism groups over $\F_p$ and over $\F_{p^2}$. The decoding
results of Sections~\ref{sec:RS} and~\ref{sec:fp2-codes} depend on the
orders of these groups and of the stabilizers of $P_\infty$.

Throughout, $p\ge5$ is prime, $g=(p-1)/2\ge2$, and $X$ is the smooth
projective model of \eqref{eq:curve}.  (For $p=3$ the curve has genus one
and a different automorphism group.)  Since $x^p-x$ is squarefree, $X$ is
hyperelliptic of genus $g$.  The hyperelliptic map $x:X\to\PP^1$ has
branch locus $\PP^1(\F_p)$, and a unique point $P_\infty$ lies over
$\infty$.  The finite branch points are $P_a=(a,0)$ for $a\in\F_p$.

\begin{remark}[Point labels in \cite{Jo05}]\label{rem:original-notation}
The note \cite{Jo05} lists
$X(\F_p)=\{(1:0:0),(0:0:1),(1:0:1),\dots,(p-1:0:1)\}$, so $(a:0:1)$
denotes the affine point $P_a$ and $(1:0:0)$ denotes the point at
infinity.  It then puts $P_1=(1:0:1)$ and states that the stabilizer of
$P_1$ is generated by $(x,y)\mapsto(x,-y)$, $(a^2x,ay)$ and $(x+1,y)$.
The translation $x\mapsto x+1$ does not fix $(1,0)$, but these maps do
generate the stabilizer of $P_\infty$.  We therefore read $P_1$ in
\cite{Jo05} as $P_\infty$ throughout, and likewise read the point
$(1:0:0)$ in conjecture (C2).  (In the plane projective closure of
\eqref{eq:curve} the point at infinity is $[0:1:0]$; divisor computations
use the smooth model.)
\end{remark}

\begin{lemma}\label{lem:poles}
$\ord_{P_\infty}(x)=-2$ and $\ord_{P_\infty}(y)=-p$.  For every $m\ge0$
the set
\begin{equation}\label{eq:monomial-basis}
\{x^i:\ 2i\le m\}\ \cup\ \{x^iy:\ 2i+p\le m\}
\end{equation}
is an $F$-basis of $L(mP_\infty)$.  In particular $\ell(mP_\infty)$ is the
number of non-gaps $\le m$ of the numerical semigroup
$\langle2,p\rangle$.
\end{lemma}
\begin{proof}
The map $x$ has degree two and is totally ramified at $P_\infty$, so
$\ord_{P_\infty}(x)=-2$.  The equation $y^2=x^p-x$ then gives
$\ord_{P_\infty}(y)=-p$.  The affine curve $y^2=x^p-x$ is smooth,
because $x^p-x$ is squarefree and the characteristic is odd.  Hence its
coordinate ring $F[x,y]/(y^2-x^p+x)$ is integrally closed and equals the
ring of functions regular away from $P_\infty$.  Each element has a unique
expression $A(x)+yB(x)$.  The monomials in \eqref{eq:monomial-basis}
have distinct pole orders (even for $x^i$, odd for $x^iy$), so leading
terms cannot cancel.  It follows that $A+yB\in L(mP_\infty)$ exactly when
it is a combination of these monomials.  Compare
\cite[Prop.~VI.4.1]{St93}.
\end{proof}

\begin{definition}
A finite group $G\le\Aut_F(X)$ is \emph{large} (over $F$) if
$|G|>|X(F)|$.
\end{definition}

\begin{lemma}\label{lem:large}
If $G$ is large, every $F$-rational point is ramified in $X\to X/G$.
\end{lemma}
\begin{proof}
A point $P\in X(F)$ with trivial stabilizer would have an orbit of size
$|G|>|X(F)|$, which is impossible.  A nontrivial stabilizer is the
inertia group of $P$ in the Galois cover $X\to X/G$.
\end{proof}

\subsection{\texorpdfstring{Automorphisms over $\F_p$}{Automorphisms over F\_p}}\label{subsec:fp}

This subsection determines the automorphism group $G_p$ of $X$ over
$\F_p$ and the stabilizer $H_p$ of $P_\infty$. The main point is a
criterion for when a M\"obius transformation of the $x$-line lifts to
$X$.

Let $F=\F_p$.  Then $X(\F_p)=\{P_\infty\}\cup\{P_a:a\in\F_p\}$ has
$p+1$ points, because $x^p-x=0$ for $x\in\F_p$.  For
$M=\begin{pmatrix}a&b\\c&d\end{pmatrix}\in\GL_2(\F_p)$ put
$x'=(ax+b)/(cx+d)$.  Since the entries lie in $\F_p$,
\begin{equation}\label{eq:mobius-identity}
x'^{\,p}-x'=\frac{(ax^p+b)(cx+d)-(ax+b)(cx^p+d)}{(cx+d)^{p+1}}
=\frac{(\det M)\,(x^p-x)}{(cx+d)^{p+1}}.
\end{equation}
Thus any lift of $x\mapsto x'$ to $X$ has the form
\begin{equation}\label{eq:lift}
(x,y)\longmapsto\Bigl(\frac{ax+b}{cx+d},\ \frac{\eta\,y}{(cx+d)^{(p+1)/2}}\Bigr),
\qquad\eta^2=\det M.
\end{equation}
Rescaling $M$ by $\lambda\in\F_p^\times$ multiplies $\det M$ by
$\lambda^2$, so the square class of $\det M$ is well defined in
$\PGL_2(p)$.  Hence \emph{the projectivity lifts to an
$\F_p$-automorphism if and only if $\det M$ is a square in $\F_p^\times$},
that is, if and only if its class lies in $\PSL_2(p)$.  Each such class
has exactly two lifts, which differ by $\iota$.

Conversely, every automorphism of $X$ (over $\overline{\F}_p$) normalizes
$\iota$ and induces a M\"obius transformation preserving the branch locus
$\PP^1(\F_p)$.  A M\"obius transformation is determined by the images of
$0,1,\infty$, so it lies in $\PGL_2(p)$.  Therefore the group
$G_p=\Aut_{\F_p}(X)$ fits into an exact sequence
\[
1\to\langle\iota\rangle\to G_p\to\PSL_2(p)\to1,
\qquad |G_p|=p(p^2-1),
\]
in agreement with \cite{BrS86,Go03}.  Explicit elements are
\begin{equation}\label{eq:generators}
\begin{gathered}
\iota:(x,y)\mapsto(x,-y),\qquad
\delta_a:(x,y)\mapsto(a^2x,ay),\qquad
\tau_b:(x,y)\mapsto(x+b,y),\\
\omega:(x,y)\mapsto\bigl(-1/x,\ y/x^{(p+1)/2}\bigr),
\end{gathered}
\end{equation}
for $a\in\F_p^\times$ and $b\in\F_p$.  Note that $\delta_{-1}=\iota$.
The group $G_p$ is transitive on $X(\F_p)$.

The stabilizer $H_p=\Stab_{G_p}(P_\infty)$ is generated by the
$\delta_a$ and $\tau_b$ and has order
\begin{equation}\label{eq:Hp-order}
|H_p|=|G_p|/(p+1)=p(p-1).
\end{equation}
(The value $2p(p-1)$ given in \cite{Jo05} is the order of the
corresponding stabilizer over $\F_{p^2}$; see
Section~\ref{subsec:fp2}.)  Since $\iota$ fixes every branch point, the
group of permutations that $H_p$ induces on $\{P_a\}$ is
\begin{equation}\label{eq:square-affine}
\{x\mapsto ux+b:\ u\in(\F_p^\times)^2,\ b\in\F_p\},
\qquad\text{of order }p(p-1)/2.
\end{equation}
A nonsquare multiplier $u$ cannot be lifted to an $\F_p$-automorphism
fixing $P_\infty$, by the criterion above.

\begin{proposition}\label{prop:H-comp}
For every $m\ge0$, the $H_p$-module $L(mP_\infty)$ has a composition
series with one-dimensional factors.  Hence its semisimplification is a
direct sum of characters of $H_p$.
\end{proposition}
\begin{proof}
Order the basis \eqref{eq:monomial-basis} by increasing pole order at
$P_\infty$.  Since $H_p$ fixes $P_\infty$, each subspace of functions with
pole order at most $j$ is $H_p$-stable.  The pole orders are distinct, so
the successive quotients are one-dimensional.  Explicitly, with
$\epsilon\in\{0,1\}$,
\[
\delta_a\cdot(x^iy^\epsilon)=a^{-2i-\epsilon}x^iy^\epsilon,\qquad
\tau_b\cdot(x^iy^\epsilon)=(x-b)^iy^\epsilon,\qquad
\iota\cdot(x^iy^\epsilon)=(-1)^\epsilon x^iy^\epsilon,
\]
using the convention \eqref{eq:convention}.  Note that $\iota$ acts
nontrivially on the basis elements involving $y$; the corresponding
statement in \cite{Jo05} is incorrect.
\end{proof}

\subsection{\texorpdfstring{Automorphisms and rational points over $\F_{p^2}$}{Automorphisms and rational points over F\_\{p\textasciicircum 2\}}}
\label{subsec:fp2}

Over $\F_{p^2}$ every M\"obius transformation in $\PGL_2(p)$ lifts, so
the automorphism group is twice as large. This subsection finds that
group, counts the rational points and describes the orbits. The regular
orbit found here carries the codes of Section~\ref{sec:fp2-codes}.

Let $F=\F_{p^2}$.  Every element of $\F_p^\times$ is a square in $F$, so
every class in $\PGL_2(p)$ lifts by \eqref{eq:lift}, and by the argument
of Section~\ref{subsec:fp} these are all the automorphisms:
\begin{equation}\label{eq:gq-order}
1\to\langle\iota\rangle\to G_{p^2}\to\PGL_2(p)\to1,
\qquad |G_{p^2}|=2p(p^2-1).
\end{equation}
The stabilizer $H=\Stab_{G_{p^2}}(P_\infty)$ is the preimage of the Borel
subgroup and has order $2p(p-1)$.  It consists of the maps
\begin{equation}\label{eq:H-affine}
(x,y)\longmapsto(ux+b,\ \eta y),
\qquad u\in\F_p^\times,\ b\in\F_p,\ \eta\in\F_{p^2},\ \eta^2=u.
\end{equation}

\begin{proposition}\label{prop:orbits}
Suppose $p\equiv3\pmod4$.  The $G_{p^2}$-orbits on $X(\F_{p^2})$ are
\[
O_1=X(\F_p),\quad|O_1|=p+1;\qquad
O_2=X(\F_{p^2})\setminus X(\F_p),\quad|O_2|=2p(p-1).
\]
The subgroup $H$ acts regularly on $O_2$.  Consequently the stabilizer
$H_2$ of a point $P_2\in O_2$ has order $p+1$, $H\cap H_2=\{1\}$, and
$G_{p^2}=H\cdot H_2$.  If $p\equiv1\pmod4$, then $X(\F_{p^2})=X(\F_p)$ is
a single orbit of size $p+1$.
\end{proposition}
\begin{proof}
For $x\in\F_p$ we get $y=0$.  For $x\in\F_{p^2}\setminus\F_p$, put
$z=x^p-x\ne0$.  Then $z^p=x-x^p=-z$, so $z^{p-1}=-1$ and
\[
z^{(p^2-1)/2}=(z^{p-1})^{(p+1)/2}=(-1)^{(p+1)/2}.
\]
Hence $z$ is a nonzero square in $\F_{p^2}$ exactly when $p\equiv3\pmod4$.
In that case each of the $p^2-p$ values $x\notin\F_p$ has two points
above it, and otherwise none.

Write $x\notin\F_p$ uniquely as $\alpha+\beta\theta$, with
$\alpha\in\F_p$, $\beta\in\F_p^\times$, for a fixed
$\theta\in\F_{p^2}\setminus\F_p$.  The maps $x\mapsto ux+b$ with
$u\in\F_p^\times$, $b\in\F_p$ send $(\alpha,\beta)$ to
$(u\alpha+b,u\beta)$, so they act regularly on these $p(p-1)$ values.
For a fixed affine map the two lifts in \eqref{eq:H-affine}, with
$\pm\eta$, send a given point to the two points over the image value.
Hence $H$ is transitive on $O_2$, and since $|H|=|O_2|$ it acts
regularly.  The set $X(\F_p)$ is the fixed locus of $\iota$ together with
$P_\infty$, so it is $G_{p^2}$-stable, and it is a single orbit by
Section~\ref{subsec:fp}.  The remaining assertions follow from
orbit--stabilizer.
\end{proof}

\begin{remark}
This agrees with the Hasse--Weil bound.  For $p\equiv3\pmod4$,
$|X(\F_{p^2})|=2p^2-p+1=p^2+1+2g\,p$, so $X$ is \emph{maximal} over
$\F_{p^2}$.  For $p\equiv1\pmod4$, $|X(\F_{p^2})|=p+1=p^2+1-2g\,p$, so
$X$ is \emph{minimal}.  The orbit statement was originally communicated
by R.~Guralnick \cite{Gu04}; the proof above is self-contained.
\end{remark}

For $p\equiv3\pmod4$, the invariant divisor $D_2=\sum_{P\in O_2}P$ has
degree $2p(p-1)>2g-2$, so $\ell(D_2)=2p(p-1)+1-g$.

\begin{example}[The case $p=7$]\label{ex:curve-p7}
Let $p=7$, so $g=3$. The semigroup $\langle2,7\rangle$ has gaps $1,3,5$.
By Lemma~\ref{lem:poles}, $L(7P_\infty)$ has basis $\{1,x,x^2,x^3,y\}$
and $\ell(7P_\infty)=5$. This agrees with Riemann--Roch, since
$7>2g-2=4$ and $7+1-3=5$.

Over $\F_7$ the curve has $8$ rational points and $|G_7|=7\cdot48=336$.
Since $336>8$, the group $G_7$ is large. The stabilizer $H_7$ has order
$42$, and it induces $21$ permutations of the points $P_a$, by
\eqref{eq:square-affine}.

Over $\F_{49}$ the curve has $92$ rational points and $|G_{49}|=672$, so
$G_{49}$ is large as well. The orbits have sizes $|O_1|=8$ and
$|O_2|=84$. The stabilizer $H$ has order $84$ and acts regularly on
$O_2$. The divisor $D_2$ has degree $84$, and $\ell(D_2)=84+1-3=82$.
\end{example}

%======================================================================
\section{AG codes and group representations}\label{sec:modules}

This section relates three actions of a group of automorphisms: on the
evaluation points, on the code and on the Riemann--Roch space. It then
shows by example that the last of these need not be a permutation
representation.

Let $G\le\Aut_F(X)$ be finite, and let $D$ and $E=\sum_{P\in\Omega}P$ be
$G$-invariant with disjoint supports, where $\Omega$ is a set of rational
points.  For $g\in G$ let $\pi_E(g)\in S_\Omega$ be the coordinate
permutation $(\pi_E(g)c)_P=c_{g^{-1}P}$.  With \eqref{eq:convention},
\begin{equation}\label{eq:equivariance}
\ev_E(g\cdot f)=\pi_E(g)\,\ev_E(f)\qquad(f\in L(D)).
\end{equation}
Hence $\pi_E(G)\le\PAut(C_X(D,E))$.

\begin{proposition}\label{prop:actions}
Write $C=C_X(D,E)$, and let $\rho_C:G\to\GL(C)$ be the restriction of
$\pi_E$ to $C$.
\begin{enumerate}
\item $\ker\pi_E$ is the pointwise stabilizer of $\Omega$.  This holds by
the definition of a permutation action and needs no hypothesis on $L(D)$.
\item $\ker\pi_E\le\ker\rho_C$, with equality if $L(D)$ separates the
points of $\Omega$.
\item If $\ev_E$ is injective (for example if $\deg D<n$), then $\ev_E$
is an isomorphism of $F[G]$-modules $L(D)\cong C$.  Since $C$ is a
submodule of the permutation module $F^\Omega$, so is $L(D)$, up to
isomorphism.
\end{enumerate}
\end{proposition}
\begin{proof}
(1) is the definition.  For (2), if $g\in\ker\rho_C$ then
$f(g^{-1}P)=f(P)$ for all $f\in L(D)$ and $P\in\Omega$, and separation
forces $g^{-1}P=P$.  (3) follows from \eqref{eq:equivariance}.
\end{proof}

\begin{remark}\label{rem:separation}
The statement in \cite{Jo05} that ``$D$ ample implies $L(D)$ separates
points'' requires correction.  On a curve every divisor of positive
degree is ample, but its complete linear system need not separate points.
For example, if $g\ge1$ then $L(P_\infty)=F$ separates no pair of points.
A standard sufficient condition is $\deg D\ge2g+1$, which makes $D$ very
ample \cite[Cor.~IV.3.2]{Ha77}.  It is also known that a nontrivial
automorphism of a curve of genus $g\ge2$ has at most $2g+2$ fixed points,
so $\pi_E$ is faithful when $n>2g+2$ \cite[\S VII.3]{St93}.
\end{remark}

\subsection{Invariant submodules need not be permutation modules}
\label{subsec:not-perm}

This subsection corrects a claim of~\cite{Jo05} about the representation
on $L(D)$ and gives a counterexample with $p=7$.

The note \cite{Jo05} concluded, from the result of Wesemeyer \cite{We98}
under hypotheses such as $n>\max(2\deg D,2g+2)$, that the representation
of $G$ on $L(D)$ is equivalent to one by permutation matrices.
Proposition~\ref{prop:actions}(3) only gives a \emph{submodule} of a
permutation module, and the stronger conclusion is false.

\begin{example}[A Riemann--Roch module that is not a permutation module]
\label{ex:not-permutation}
Let $p=7$, $F=\F_{49}$, $g=3$, and let $O_1,O_2$ be as in
Proposition~\ref{prop:orbits}.  Put $D=\sum_{P\in O_1}P$ and
$E=\sum_{P\in O_2}P$.  These are $G_{49}$-invariant with disjoint
supports, and
\[
\deg D=8,\qquad n=84>\max\{16,8\}.
\]
Riemann--Roch gives $\ell(D)=8+1-3=6$, and $\ev_E$ is injective.  We have
$\operatorname{div}(y)=\sum_aP_a-7P_\infty$ and
$\operatorname{div}(x)=2P_0-2P_\infty$, so
\[
\operatorname{div}(1/y)+D=8P_\infty\ge0,\qquad
\operatorname{div}(x/y)+D=2P_0+6P_\infty\ge0.
\]
Hence $1/y$ and $x/y$ lie in $L(D)$.  The automorphism $\tau_1$ has
order $7$, and by \eqref{eq:convention}
\[
\tau_1\cdot(x/y)=\frac{x-1}{y}=\frac xy-\frac1y\ne\frac xy.
\]
So $\tau_1$ acts nontrivially on the six-dimensional space $L(D)$.  A
permutation of six letters whose order divides $7$ is the identity.  So
no representation of $G_{49}$ by $6\times6$ permutation matrices can be
equivalent to the representation on $L(D)$, although all the numerical
hypotheses of \cite{Jo05} hold.
\end{example}

\begin{remark}
If $\ev_E$ is not injective, $C$ is the quotient $L(D)/\ker\ev_E$.  The
decomposition of these modular representations into indecomposables,
when $p$ divides $|G|$, remains an interesting problem; see
\cite{JK03,JT02,Sc90}.
\end{remark}

%======================================================================
\section{Permutation decoding and a counting obstruction}\label{sec:PD}

This section defines permutation-decoding sets, describes the decoding
algorithm and proves a counting bound that every such set must satisfy.
The bound is the source of the negative results in
Sections~\ref{sec:RS} and~\ref{sec:fp2-codes}.

Let $C$ be an $[n,k,d]$ code with $t=\lfloor(d-1)/2\rfloor\ge1$, an
information set $\cI$ and check set $\cJ$ ($|\cJ|=r=n-k$), and systematic
matrices $\Gsys=(I_k\mid A)$, $\Hsys=(-A^{\mathsf T}\mid I_r)$.

\begin{lemma}[Syndrome criterion]\label{lem:syndrome}
Let $v=c+e$ with $c\in C$ and $\wt(e)\le t$.  The information coordinates
of $v$ agree with those of $c$ if and only if
$\wt(\Hsys v^{\mathsf T})\le t$.
\end{lemma}
\begin{proof}
Write $e=(u,e_\cJ)$ and $s=\Hsys v^{\mathsf T}=\Hsys e^{\mathsf T}
=e_\cJ-A^{\mathsf T}u$.  If $u=0$ then $s=e_\cJ$ has weight at most $t$.
Conversely, $(u,A^{\mathsf T}u)\in C$ and $(u,A^{\mathsf T}u)=e-(0,s)$.
This codeword has weight at most $\wt(e)+\wt(s)\le2t<d$, so it is zero
and $u=0$.
\end{proof}

\begin{definition}\label{def:pd-set}
For $1\le s\le t$, a subset $S\subseteq\PAut(C)$ is an
\emph{$s$-PD-set} relative to $\cI$ if for every set $T$ of at most $s$
coordinates some $\sigma\in S$ satisfies $\sigma(T)\subseteq\cJ$.  A
$t$-PD-set is a \emph{full PD-set}.
\end{definition}

\paragraph{Algorithm.}
Given $v$ and an $s$-PD-set $S$, compute $\wt(\Hsys(\sigma v)^{\mathsf
T})$ for $\sigma\in S$ until a value $\le t$ is found.  Then re-encode the
information symbols of $\sigma v$ to get $c_\sigma\in C$, and output
$\sigma^{-1}c_\sigma$.  By Lemma~\ref{lem:syndrome} this decodes
correctly whenever $\wt(e)\le s$.  (Any $\sigma$ that passes the test
gives the correct answer, since $\wt(\sigma e)=\wt(e)\le t$.)

\begin{theorem}[Counting bound]\label{thm:counting}
If $S$ is an $s$-PD-set relative to $\cI$, then
\begin{equation}\label{eq:counting}
|S|\binom{n-k}{s}\ \ge\ \binom ns.
\end{equation}
Hence if $K\le\PAut(C)$ satisfies $|K|\binom{n-k}s<\binom ns$, then no
subset of $K$ is an $s$-PD-set for any information set.
\end{theorem}
\begin{proof}
A fixed $\sigma$ sends exactly $\binom rs$ of the $s$-subsets into $\cJ$,
namely the $s$-subsets of $\sigma^{-1}(\cJ)$.  All $\binom ns$ subsets
must be covered.
\end{proof}

\begin{remark}
Inequality \eqref{eq:counting} is the simplest case of Gordon's covering
argument \cite{Gor82}.  The iterated (Sch\"onheim-type) version, used for
example in \cite{KMM05}, is sharper.  It is only a \emph{necessary}
condition; Example~\ref{ex:p17} shows that it is not sufficient.
\end{remark}

\begin{example}[A $2$-PD-set of seven translations]\label{ex:pd-translations}
Let $C=\RS(\F_7,3)$, with coordinates indexed by $\F_7$ as in
Example~\ref{ex:rs7}. It is a $[7,3,5]$ MDS code, so $t=2$ and every
$3$-subset is an information set. Take $\cI=\{0,1,2\}$ and
$\cJ=\{3,4,5,6\}$. Let $S$ be the set of the seven translations
$x\mapsto x+b$, which lie in $\PAut(C)$ by
Section~\ref{subsec:bg-rs}. A pair $T=\{a_1,a_2\}$ can be translated
into $\cJ$ exactly when $\cJ$ contains two elements whose difference is
$\pm(a_1-a_2)$. The differences $4-3=1$, $5-3=2$ and $6-3=3$ represent
every nonzero class of $\F_7$ up to sign. A single position can be
translated to $3$. So $S$ is a $2$-PD-set relative to $\cI$, and hence a
full PD-set. The counting bound holds with room to spare:
$7\binom42=42\ge21=\binom72$. The same bound shows that no set of three
permutations is a $2$-PD-set for this code, since $3\binom42=18<21$.
\end{example}

%======================================================================
\section{\texorpdfstring{Codes over $\F_p$}{Codes over F\_p}}\label{sec:RS}

This section treats the codes of conjecture (C1). It identifies them as
Reed--Solomon codes and proves that two errors can always be corrected by
permutation decoding. It then reports exhaustive computations for
$p\le23$ and ends with a counterexample to (C1) at $p=37$.

Let $E=\sum_{a\in\F_p}P_a$ and $D=mP_\infty$ with $m\ge0$.

\begin{theorem}\label{thm:RS}
With $k=\min\bigl(p,\lfloor m/2\rfloor+1\bigr)$,
\[
C_X(mP_\infty,E)=\RS(\F_p,k),
\]
a $[p,k,p-k+1]$ code.  For $2\le k\le p-2$ its full permutation group is
$\AGL(1,p)$, of order $p(p-1)$.
\end{theorem}
\begin{proof}
Every basis monomial in \eqref{eq:monomial-basis} involving $y$ vanishes
at the points $P_a$.  So the code is spanned by the evaluations of
$1,x,\dots,x^{\lfloor m/2\rfloor}$ on $\F_p$.  Since $x^p\equiv x$ on
$\F_p$, its dimension is $k$.  The permutation group is the one-variable
case of the theorem of Kn\"orr and Willems on generalized Reed--Muller
codes over prime fields \cite{KW90}, whose exceptional cases correspond to
$k\in\{1,p-1,p\}$; see also \cite{CLLS25,BC93}.
\end{proof}

\begin{remark}\label{rem:RSkernel}
When $m\ge p$, the function $y\in L(mP_\infty)$ evaluates to zero on
$E$, so $\ev_E$ is not injective.  The code group $\AGL(1,p)$ is twice
as large as the geometric image \eqref{eq:square-affine}.  So the full
code symmetry is \emph{not} induced by $\Aut_{\F_p}(X)$, and the curve
structure adds nothing beyond the Reed--Solomon structure for these codes.
\end{remark}

\begin{theorem}[Two errors]\label{thm:two-pd}
Let $1\le k\le p-4$.  Then $\AGL(1,p)\le\PAut(\RS(\F_p,k))$ is a
$2$-PD-set relative to every information set.  More precisely, for each
set $T$ of two coordinates exactly $r(r-1)$ affine maps send $T$ into
$\cJ$, where $r=p-k$.
\end{theorem}
\begin{proof}
Since $d=p-k+1\ge5$, we have $t\ge2$.  Every $k$-set is an information
set because the code is MDS.  Given $T=\{a_1,a_2\}$ and any ordered pair
$(b_1,b_2)$ of distinct check positions, sharp $2$-transitivity gives a
unique affine map with $a_i\mapsto b_i$.  There are $r(r-1)$ such pairs.
Single errors are handled by transitivity.
\end{proof}

\begin{example}[{The $[7,3,5]$ code of \cite{Jo05}}]\label{ex:7-3-5}
For $p=7$ and $D=4P_\infty$, $L(D)=\langle1,x,x^2\rangle$ and
$C=\RS(\F_7,3)$.  The systematic matrix computed in GUAVA \cite{Jo05} is
\[
\Gsys=\begin{pmatrix}
1&0&0&6&4&6&4\\
0&1&0&6&3&3&1\\
0&0&1&3&1&6&3
\end{pmatrix}.
\]
We checked by direct computation that this code has minimum weight $5$,
that the permutations $(1,5,4)(2,3,7)$ and $(1,3,2,6,7,4)$ preserve it
and generate a group of order $42$, and that a search of all $5040$
elements of $S_7$ finds no further permutation automorphisms.  Thus
$\PAut(C)\cong\AGL(1,7)$.  Here $t=2$, so by Theorem~\ref{thm:two-pd} the
whole group is a full PD-set, and $4\cdot3=12$ of its elements handle
each pair of error positions.  (The note \cite{Jo05} proposed
$\PAut(C)\setminus\{1\}$, which also works.)
\end{example}

\subsection{\texorpdfstring{Exhaustive computations for small $p$}{Exhaustive computations for small p}}

This subsection reports a computer search. For each prime $p\le23$ it
decides which information sets make $\AGL(1,p)$ a full PD-set.

For $\RS(\F_p,k)$, every $k$-subset of $\F_p$ is an information set, and
whether $\AGL(1,p)$ is a full PD-set depends only on the
$\AGL(1,p)$-orbit of the information set.  We tested every orbit.

\begin{proposition}[computer verified]\label{prop:small-p}
Let $5\le p\le23$ be prime and $2\le k\le p-2$, and let
$t=\lfloor(p-k)/2\rfloor$.
\begin{enumerate}
\item For every such $(p,k)$ there is an information set relative to
which $\AGL(1,p)$ is a full PD-set for $\RS(\F_p,k)$.
\item For $p\le13$, every information set has this property.
\item For $p\in\{17,19,23\}$, the information sets for which $\AGL(1,p)$
fails to be a full PD-set occur exactly for the following
$(p,k)$:
\[
\begin{array}{cccc}
\toprule
(p,k) & t & \text{failing orbits / all orbits} &
\text{failing information sets / all}\\
\midrule
(17,7) & 5 & 3/75 & 680/19448\\
(17,9) & 4 & 5/95 & 1122/24310\\
(19,7) & 6 & 13/154 & 3591/50388\\
(19,9) & 5 & 22/280 & 5510/92378\\
(19,10) & 4 & 2/280 & 209/92378\\
(19,11) & 4 & 35/228 & 10944/75582\\
(19,13) & 3 & 3/86 & 570/27132\\
\midrule
(23,7) & 8 & 11/492 & 4807/245157\\
(23,8) & 7 & 14/984 & 6072/490314\\
(23,9) & 7 & 484/1630 & 237567/817190\\
(23,10) & 6 & 187/2282 & 88803/1144066\\
(23,11) & 6 & 2141/2694 & 1072260/1352078\\
(23,12) & 5 & 345/2694 & 168291/1352078\\
(23,13) & 5 & 1707/2282 & 853116/1144066\\
(23,14) & 4 & 42/1630 & 19228/817190\\
(23,15) & 4 & 396/984 & 193798/490314\\
(23,17) & 3 & 5/207 & 2024/100947\\
\bottomrule
\end{array}
\]
\end{enumerate}
In each listed case the counting bound \eqref{eq:counting} is satisfied.
\end{proposition}

For $p\le19$ the computation was done independently by two of the
authors.  The full search uses two reductions: it suffices to test one
information set per $\AGL(1,p)$-orbit, and, for a fixed information set,
one error pattern $T$ per $\AGL(1,p)$-orbit (if $\sigma(T)\subseteq\cJ$
then $(\sigma\rho^{-1})(\rho(T))\subseteq\cJ$).  With these, $p=23$
takes under a minute in Python.  The code and complete output are
available in the repository \cite{Jo26}; a short check of
Example~\ref{ex:p17} is in Appendix~\ref{app:script}.

At $p=23$, failure has become common: for $k=11$ about $79\%$ of all
information sets fail, although some information set still works for
every $k$.

\begin{remark}\label{rem:interval}
The most natural choice, the ``interval'' information set
$\cI=\{0,1,\dots,k-1\}$, works for every $k$ when $p\le17$, but fails
for $k\in\{7,9,11\}$ when $p=19$ and for
$k\in\{7,9,10,\dots,15\}$ when $p=23$.  So a good information set must
be chosen with some care even for moderate $p$.
\end{remark}

\begin{example}[Failure despite the counting bound]\label{ex:p17}
Let $p=17$, $k=9$, so $C=\RS(\F_{17},9)$ is a $[17,9,9]$ code with
$t=4$.  Take check positions
\[
\cJ=\{0,1,2,3,4,7,9,10\}\subset\F_{17},
\]
so that $\cI=\F_{17}\setminus\cJ$.  An enumeration of all $272$ affine
maps shows that none sends the error support $T=\{0,1,4,5\}$ into $\cJ$.
Thus $\AGL(1,17)$ is not a $4$-PD-set relative to $\cI$, although
\[
272\binom84=19040>2380=\binom{17}4.
\]
\end{example}

\subsection{A counterexample to (C1) for every information set}

This subsection shows that (C1) fails. At $p=37$ the counting bound alone
rules out every information set.

\begin{theorem}\label{thm:counter-RS}
Let $p=37$ and $D=32P_\infty$, so $C=\RS(\F_{37},17)$ with parameters
$[37,17,21]$, $r=20$ and $t=10$.  No subset of $\PAut(C)=\AGL(1,37)$ is a
full PD-set, relative to any information set.  Hence conjecture (C1) is
false.
\end{theorem}
\begin{proof}
$|\AGL(1,37)|\binom{20}{10}=1332\cdot184756=246094992
<348330136=\binom{37}{10}$.  Apply Theorem~\ref{thm:counting}.
\end{proof}

\begin{remark}
A direct check shows that $p=37$ is the smallest prime for which the
counting bound, applied to $\AGL(1,p)$ with $s=t$, rules out a full
PD-set in some dimension $2\le k\le p-2$.  At $p=37$ this happens exactly
for $k\in\{17,\dots,25\}\cup\{27\}$, so $k=17$ is the smallest case.  For
$29\le p\le31$ the counting bound decides nothing, and the exhaustive
check is beyond a simple Python program; see \cite{Jo26} for what a
compiled implementation needs.
\end{remark}

%======================================================================
\section{\texorpdfstring{Codes over $\F_{p^2}$}{Codes over F\_\{p\textasciicircum 2\}}}\label{sec:fp2-codes}

This section treats the codes of conjecture (C2), which live on the
regular orbit $O_2$ of Proposition~\ref{prop:orbits}. The main result is
that a regular group of permutations can correct only about $\log n$
errors by permutation decoding. This gives an infinite family of
counterexamples to (C2).

Assume $p\equiv3\pmod4$ and $p\ge7$.  Put $E_2=\sum_{P\in O_2}P$, so
$n=2p(p-1)$, and let $D=mP_\infty$ with $m<n$.  The group $H$ of
\eqref{eq:H-affine} fixes $D$ and $E_2$.  By
Proposition~\ref{prop:orbits} it acts regularly on the $n$ coordinates,
so its image $\pi_{E_2}(H)$ has exactly $n$ elements.  The obstruction
below applies to any regular group.

\begin{theorem}[Regular groups correct few errors]\label{thm:regular}
Let $C$ be an $[n,k,d]$ code with $k<n$, and let $K\le\PAut(C)$.  If
$1\le s\le\min(t,n-k)$ and
\[
\Bigl(\frac{n}{n-k}\Bigr)^{s}>|K|,
\]
then no subset of $K$ is an $s$-PD-set, relative to any information set.
In particular, if $|K|=n$ this holds for every
$s>\log n/\log\bigl(n/(n-k)\bigr)$ with $s\le\min(t,n-k)$.
\end{theorem}
\begin{proof}
For $0\le i<s\le n-k$ we have $(n-i)/(n-k-i)\ge n/(n-k)$, so
\[
\binom ns\Big/\binom{n-k}s=\prod_{i=0}^{s-1}\frac{n-i}{n-k-i}
\ge\Bigl(\frac n{n-k}\Bigr)^s>|K|.
\]
Apply Theorem~\ref{thm:counting}.
\end{proof}

\begin{corollary}[An infinite family of counterexamples to (C2)]
\label{cor:family}
Let $p\equiv3\pmod4$, $p\ge7$, $n=2p(p-1)$, $g=(p-1)/2$ and
\[
m=\tfrac n2+g-1,\qquad C=C_X(mP_\infty,E_2).
\]
Then $C$ is an $[n,n/2,\ge n/2-g+1]$ code, with
$t\ge\lfloor(p-1)(2p-1)/4\rfloor$.  For
$s_0=\lfloor\log_2n\rfloor+1$ we have $s_0\le t$, and $\pi_{E_2}(H)$ is
not an $s_0$-PD-set, hence not a full PD-set, relative to any information
set.
\end{corollary}
\begin{proof}
Since $2g-1\le m<n$, the map $\ev_{E_2}$ is injective and
$k=\ell(D)=m+1-g=n/2$.  Also $d\ge n-m=n/2-g+1$, and
$n/2-g=(p-1)(2p-1)/2$, which gives the bound on $t$.  Next,
$s_0\le\log_2(2p^2)+1=2+2\log_2p$.  The inequality
$3+2\log_2p\le(p-1)(2p-1)/4$ holds at $p=7$ (where $8.62\le19.5$), and
its right side grows faster for $p\ge7$.  Since
$t>(p-1)(2p-1)/4-1$, this gives $s_0\le t$.  Finally
$(n/(n-k))^{s_0}=2^{s_0}>n=|\pi_{E_2}(H)|$, and
Theorem~\ref{thm:regular} applies.
\end{proof}

\[
\begin{array}{rrrrr}
\toprule
p & n & k & \text{lower bound for }t & s_0\\
\midrule
7 & 84 & 42 & 19 & 7\\
11 & 220 & 110 & 52 & 8\\
19 & 684 & 342 & 166 & 10\\
23 & 1012 & 506 & 247 & 10\\
31 & 1860 & 930 & 457 & 11\\
\bottomrule
\end{array}
\]

\begin{example}[The case $p=7$]\label{ex:p7}
For $p=7$, Corollary~\ref{cor:family} gives $m=44$ and the
$[84,42,\ge40]$ code $C_X(44P_\infty,E_2)$ over $\F_{49}$, with $t\ge19$.
Here
\[
84\binom{42}{7}=2266179552<4529365776=\binom{84}{7},
\]
while $s=6$ is not excluded by the counting bound.  So the stabilizer
image is not a $7$-PD-set.
\end{example}

\begin{remark}
These results concern the geometric stabilizer proposed in (C2), not the
full group $\PAut(C)$, which we have not determined.  The stabilizer of
$D=mP_\infty$ in $G_{p^2}$ is exactly $H$, so if $\PAut(C)$ is induced by
automorphisms of $X$ preserving $D$ and $E_2$, then
$\PAut(C)=\pi_{E_2}(H)$ and Corollary~\ref{cor:family} applies to the full
permutation group.  The results of \cite{We98,JK06} that give such an
identification assume $n$ large compared with $\deg D$, and this fails
for the half-rate family ($n<2\deg D$).  So the question remains open
(Problem~\ref{prob:paut}).  In any case Theorem~\ref{thm:regular} shows
that any group $K\le\PAut(C)$ of polynomial size in $n$ can support
permutation decoding of only $O(\log n)$ errors at fixed rate.
\end{remark}

%======================================================================
\section{Complexity and computation of the action}\label{sec:complexity}

This section discusses the cost of permutation decoding for the codes
above and the cost of finding their automorphisms.

The conjectures of \cite{Jo05} combined a group-size estimate with the
claim that a group is a full PD-set.  These must be separated.  If
syndromes are computed with a dense $r\times n$ matrix, one trial costs
$O(nr)$ field operations plus $O(n)$ coordinate work.  A worst-case scan
of $S$ therefore costs $O(|S|\,nr)$ field operations.  The estimates
$|\AGL(1,p)|=O(p^2)$ and $|H|=O(p^2)$ describe group orders, not total
costs, and they do not show asymptotic superiority over Reed--Solomon or
AG decoders \cite{HvLP98,GS99}.  The alphabet sizes $p$ and $p^2$ also grow
with $n$, which matters for bit complexity.

\begin{example}\label{ex:cost}
For the $[7,3,5]$ code of Example~\ref{ex:7-3-5} we have $n=7$, $r=4$ and
$|\PAut(C)|=42$. One trial costs about $nr=28$ field operations. A scan of
the whole group costs at most about $42\cdot28=1176$ field operations.
For $\RS(\F_{23},11)$ the same count gives $nr=23\cdot12=276$ and
$506\cdot276=139656$.
\end{example}

On the other hand, a general algorithm for the automorphism group of a
linear code, such as Leon's \cite{Le82}, can be expensive.  (The GUAVA
command \texttt{PermutationAutomorphismGroup} took over a day for $p=13$
on 2005 hardware \cite{Jo05}.)  For the codes considered here, the
formulas \eqref{eq:generators} and \eqref{eq:H-affine} give the action
of $G_p$, $H$ and $G_{p^2}$ on coordinates directly, at negligible cost.
This remains a practical advantage of geometric constructions.

%======================================================================
\section{Summary of results and open problems}\label{sec:problems}

This section gathers up the paper. The first subsection lists the main
results in one place. The second lists problems that remain open.

\subsection{Summary of results}

Each item below points to the precise statements. The permutation group
in item~1 is cited from the literature, and item~3 rests in part on the
computer search of Proposition~\ref{prop:small-p}. Everything else is
proved in this paper.

\begin{enumerate}
\item The $\F_p$ codes of length $p$ are full-length Reed--Solomon codes,
and their permutation group $\AGL(1,p)$ is larger than the geometric
image (Theorem~\ref{thm:RS}, Remark~\ref{rem:RSkernel}).
\item $\AGL(1,p)$ always gives $2$-error permutation decoding
(Theorem~\ref{thm:two-pd}).
\item Full permutation decoding by $\AGL(1,p)$ depends on the information
set from $p=17$ on, though for $p\le23$ a good information set always
exists (Proposition~\ref{prop:small-p}, Example~\ref{ex:p17}), and fails for all information sets at $p=37$
(Theorem~\ref{thm:counter-RS}).
\item Over $\F_{p^2}$, regular groups such as the geometric stabilizer
correct only $O(\log n)$ errors by permutation decoding at fixed rate,
giving infinitely many counterexamples to (C2)
(Theorem~\ref{thm:regular}, Corollary~\ref{cor:family}).
\item Equivariant AG codes are submodules of permutation modules, but the
Riemann--Roch module need not be a permutation module
(Example~\ref{ex:not-permutation}).
\end{enumerate}

\subsection{Open problems}

The four problems below ask what permutation decoding can still do for
these codes. The first and third concern the codes over $\F_{p^2}$. The
second concerns the Reed--Solomon codes at the two primes not yet
decided. The fourth concerns burst errors.

\begin{problem}[Partial permutation decoding]\label{prob:partial}
For $C_X(mP_\infty,E_2)$ over $\F_{p^2}$, determine the largest
$s=s(p,m,\cI)$ for which $\pi_{E_2}(H)$ is an $s$-PD-set relative to
$\cI$.  Compare it with the bound of Theorem~\ref{thm:regular} and with
the value for the full group $\PAut(C)$, and characterize the best
information sets.
\end{problem}

\begin{problem}
For $\RS(\F_p,k)$ and $p\in\{29,31\}$, determine for which $k$ some
information set makes $\AGL(1,p)$ a full PD-set.  (By
Proposition~\ref{prop:small-p}, such a set exists for every $k$ when
$p\le23$; by Theorem~\ref{thm:counter-RS}, not when $p=37$, $k=17$.)  More generally,
describe the information sets that fail, as in
Proposition~\ref{prop:small-p} and Example~\ref{ex:p17}.
\end{problem}

\begin{problem}\label{prob:paut}
Determine $\PAut(C_X(mP_\infty,E_2))$, and decide whether it equals
$\pi_{E_2}(G_{p^2})$, for the codes of Section~\ref{sec:fp2-codes}.
\end{problem}

\begin{problem}[Burst errors]\label{prob:burst}
The curve $x^p-x=y^2$ has the separated form studied in \cite{CTO26}, and
\cite{LLMS26} constructs burst-error PD-sets on related curves.  Do the
translations $\tau_b$ and the maps $(x,y)\mapsto(ux,\eta y)$ give
PD-sets for bursts of errors supported on fibers of $x$ or $y$ in the
codes $C_X(mP_\infty,E_2)$?  How do such sets compare with the general
bound of Theorem~\ref{thm:regular}?
\end{problem}

%======================================================================
\appendix

\section{Historical GUAVA session}\label{app:guava}
The commands of \cite{Jo05}, which used GUAVA 2.0 \cite{GUAVA}, are kept
for the record, with two typos corrected: \texttt{Length(P)} is now
\texttt{Length(Pts)}, and the closing \texttt{if;} is now \texttt{fi;}.
They are version dependent, and the results of this paper do not depend
on them.
\begin{verbatim}
p:=7;; F:=GF(p);;
R:=PolynomialRing(F,2);;
indets:=IndeterminatesOfPolynomialRing(R);;
x:=indets[1];; y:=indets[2];;
Pts:=AffinePointsOnCurve(y^2-x^p+x,R,F);;
Length(Pts);
C:=OnePointAGCode(y^2-x^p+x,Pts,4,R);;
MinimumDistance(C);
G:=GeneratorMat(C);; PutStandardForm(G);; Display(G);
P:=PermutationAutomorphismGroup(C);; IdGroup(P);
Wt2:=ConstantWeightSubcode(WholeSpaceCode(p,GF(p)),2);;
for c in C do
  e:=Random(Wt2);; v:=PermutationDecode(C,c+e);;
  if c<>v then Print("\n Decoding fails. c,v = ",c," ",v,"\n"); fi;
od;
\end{verbatim}

\section{Verification script}\label{app:script}
The following Python~3 code (which also runs in Sage) checks
Example~\ref{ex:p17} and one entry of Proposition~\ref{prop:small-p}.
Information sets are reduced modulo $\AGL(1,p)$, and subsets are stored
as bit masks.  The faster program \texttt{pd\_affine.py} in
\cite{Jo26}, which also reduces error patterns modulo $\AGL(1,p)$,
produces the whole table.
\begin{verbatim}
from itertools import combinations
p, J, T = 17, {0,1,2,3,4,7,9,10}, {0,1,4,5}
print([(u,b) for u in range(1,p) for b in range(p)
       if all((u*x+b) % p in J for x in T)])     # prints []

def failing_orbits(p, k):
    maps = [[(u*x+b) % p for x in range(p)]
            for u in range(1,p) for b in range(p)]
    full, t = (1 << p) - 1, (p-k)//2
    Ts = [sum(1 << i for i in S) for S in combinations(range(p), t)]
    seen, bad, tot = set(), 0, 0
    for I in combinations(range(p), k):
        if sum(1 << i for i in I) in seen: continue
        orb = {sum(1 << m[i] for i in I) for m in maps}
        seen |= orb; tot += 1
        Js = [full ^ o for o in orb]   # images of the check set
        if not all(any(s & ~Jm == 0 for Jm in Js) for s in Ts):
            bad += 1
    return bad, tot

print(failing_orbits(17, 9))                      # (5, 95)
\end{verbatim}

%======================================================================
\section*{Glossary of technical terms}
\addcontentsline{toc}{section}{Glossary of technical terms}

This glossary lists the technical terms of the paper in alphabetical
order, each with a short explanation. Section~\ref{sec:background}
treats most of them in more detail.

\begin{description}\setlength{\itemsep}{2pt}
\item[Affine group $\AGL(1,q)$.] The group of maps $x\mapsto ux+b$ on
$\F_q$, $u\ne0$; order $q(q-1)$; sharply $2$-transitive.
\item[AG code $C_X(D,E)$.] The image of $L(D)$ under evaluation at the
points of $E$.
\item[Automorphism group $\Aut_F(X)$.] The group of automorphisms of the
curve $X$ defined over $F$.
\item[Borel subgroup.] The stabilizer of $\infty$ in $\PGL_2(p)$; the
upper-triangular classes.
\item[Branch point.] A point of $\PP^1$ over which a cover is ramified;
for $y^2=x^p-x$, the points of $\PP^1(\F_p)$.
\item[Check positions $\cJ$.] The coordinates complementary to an
information set.
\item[Counting bound.] The necessary condition
$|S|\binom{n-k}s\ge\binom ns$ for an $s$-PD-set $S$.
\item[Divisor; degree.] A formal integer combination of points; the
weighted sum of its coefficients.
\item[Equivariant (invariant) divisor.] A divisor fixed by a group of
automorphisms.
\item[Evaluation map $\ev_E$.] $f\mapsto(f(P_1),\dots,f(P_n))$.
\item[Full PD-set.] A $t$-PD-set, where $t=\lfloor(d-1)/2\rfloor$.
\item[Genus $g$.] The basic invariant of a curve; here $g=(p-1)/2$.
\item[Hamming weight $\wt(v)$.] The number of nonzero coordinates of $v$.
\item[Hasse--Weil bound.] $||X(\F_q)|-q-1|\le2g\sqrt q$.
\item[Hyperelliptic curve; involution $\iota$.] A curve of genus $\ge2$
with a degree-two map to $\PP^1$; the involution exchanging its sheets.
\item[Information set $\cI$.] $k$ coordinates on which projection of the
code is bijective.
\item[Large automorphism group.] $|G|>|X(F)|$.
\item[Maximal / minimal curve.] A curve attaining the upper / lower
Hasse--Weil bound.
\item[MDS code.] A code with $d=n-k+1$; every $k$-set is an information
set.
\item[One-point code.] An AG code with $D=mQ$ for a single point $Q$.
\item[Orbit; stabilizer.] $G\cdot\omega$ and
$G_\omega=\{g:g\omega=\omega\}$; $|G\cdot\omega|=|G|/|G_\omega|$.
\item[Permutation automorphism group $\PAut(C)$.] Coordinate
permutations preserving $C$.
\item[Permutation decoding.] Decoding by moving errors off an
information set with code automorphisms and re-encoding.
\item[Permutation module.] A module with a basis permuted by the group.
\item[$\PGL_2(p)$, $\PSL_2(p)$.] The projective general and special
linear groups over $\F_p$, of orders $p(p^2-1)$ and $p(p^2-1)/2$.
\item[Ramification.] A point of a cover is ramified if its stabilizer
(inertia group) is nontrivial.
\item[Reed--Solomon code $\RS(\F_q,k)$.] Evaluations on $\F_q$ of
polynomials of degree $<k$.
\item[Regular action.] Transitive with trivial stabilizers; then
$|G|$ equals the number of points.
\item[Riemann--Roch space $L(D)$.] Functions $f$ with
$\operatorname{div}(f)+D\ge0$, together with $0$.
\item[Semisimplification.] The direct sum of the composition factors of a
module.
\item[Separates points.] $L(D)$ separates $P\ne Q$ if some
$f\in L(D)$ has $f(P)\ne f(Q)$.
\item[Sharply $2$-transitive.] Exactly one group element maps any ordered
pair of distinct points to any other.
\item[$s$-PD-set.] A set of automorphisms moving every set of at most
$s$ coordinates into the check positions.
\item[Syndrome.] $\Hsys v^{\mathsf T}$ for a parity-check matrix $\Hsys$.
\item[Systematic form.] Generator matrix $(I_k\mid A)$ and parity-check
matrix $(-A^{\mathsf T}\mid I_{n-k})$.
\item[Very ample.] A divisor whose complete linear system embeds the
curve; on a curve, implied by $\deg D\ge2g+1$.
\item[Weierstrass point.] For a hyperelliptic curve, a fixed point of
$\iota$.
\end{description}

%======================================================================
\section*{Acknowledgments}
\addcontentsline{toc}{section}{Acknowledgments}

The first author thanks Bob Guralnick for the orbit argument behind
Proposition~\ref{prop:orbits} \cite{Gu04}, and the anonymous referee of
\cite{Jo05} for many helpful comments.  Some of the original ideas arose
in joint work with Amy Ksir and Will Traves \cite{JK06,JK03,JT02}.

This revision was prepared with the help of two AI systems, listed as
authors at the first author's request.
\emph{ChatGPT (OpenAI)} prepared the first mathematical revision.  It
identified the $\F_p$ codes as Reed--Solomon codes, introduced the
counting obstruction and the counterexamples at $p=37$ and $p=7$, proved
the $2$-PD result, gave the permutation-module counterexample and the
proof of the orbit proposition, and corrected several statements of
\cite{Jo05}.  In its response to the referee report, it found the $p=17$
information set of Example~\ref{ex:p17} and the choice
$m=n/2+g-1$ in Corollary~\ref{cor:family}.
\emph{Claude (Anthropic)} refereed that revision.  It checked the
proofs and numerical claims independently, including a brute-force
determination of the permutation group of the $[7,3,5]$ code.  It found
errors in the GAP transcript, the lift criterion, the description of
kernels, the group-action conventions, and the silent reinterpretation of
point labels.  It carried out the exhaustive computations of
Proposition~\ref{prop:small-p} (including $p=23$), proposed the
logarithmic obstruction of
Theorem~\ref{thm:regular}, verified the bibliography, located the recent
literature \cite{CTO26,LLMS26,CLLS25}, and prepared this version, including
the background section, literature review, and glossary.
Responsibility for the final content rests with the first author.

Claude (Claude Fable~5.2) also prepared Draft~2, a revision for style and
exposition that changed no statement or proof. In that draft it reordered
the sections and rewrote the front matter. It also added the examples,
orienting paragraphs and table of notation that the revision log lists as
new text in Draft~2. The script that checks the numbers in those
additions is \texttt{scripts/check\_examples\_draft2.py} in the
repository~\cite{Jo26}, at commit \texttt{9e1e1ea}.

AI review is not an independent certification of the results of this
paper. Further details on the AI usage can be found in the revision log
(\texttt{paper/revision-log.tex} in the same repository).

%======================================================================
\begin{thebibliography}{JLJHH89}\small

\bibitem[BC93]{BC93} T.~P.~Berger and P.~Charpin, The automorphism group
of generalized Reed--Muller codes, \emph{Discrete Math.} \textbf{117}
(1993), 1--17.

\bibitem[BS11]{BS11} J.~J.~Bernal and J.~J.~Sim\'on, Information sets
from defining sets in abelian codes, \emph{IEEE Trans.\ Inform.\ Theory}
\textbf{57} (2011), 7990--7999.

\bibitem[BS13]{BS13} J.~J.~Bernal and J.~J.~Sim\'on, Partial permutation
decoding for abelian codes, \emph{IEEE Trans.\ Inform.\ Theory}
\textbf{59} (2013), 5152--5170.

\bibitem[BS14]{BS14} J.~J.~Bernal and J.~J.~Sim\'on, Information sets in
abelian codes: defining sets and Groebner basis, \emph{Des.\ Codes
Cryptogr.} \textbf{70} (2014), 175--188.

\bibitem[BS23]{BS23} J.~J.~Bernal and J.~J.~Sim\'on, New advances in
permutation decoding of first-order Reed--Muller codes, \emph{Finite
Fields Appl.} \textbf{88} (2023), 102182.

\bibitem[Bi95]{Bi95} E.~Biglieri, Permutation decoding of group codes,
in \emph{Proc.\ 1995 IEEE Int.\ Symp.\ Inform.\ Theory}, p.~306.

\bibitem[BrS86]{BrS86} R.~Brandt and H.~Stichtenoth, Die
Automorphismengruppen hyperelliptischer Kurven, \emph{Manuscripta Math.}
\textbf{55} (1986), 83--92.

\bibitem[CLLS25]{CLLS25} E.~Camps-Moreno, J.~B.~Lau, H.~H.~L\'opez and
W.~Santos, The permutation group of Reed--Solomon codes over arbitrary
points, \emph{IEEE Int.\ Symp.\ Inform.\ Theory} (2025),
DOI 10.1109/ISIT63088.2025.11195458; arXiv:2601.00122.

\bibitem[CTO26]{CTO26} A.~S.~Castellanos, G.~Tizziotti and
W.~Olaya-Le\'on, Permutation decoding of AG codes from curves defined by
separated polynomials, arXiv:2608.03592 (2026).

\bibitem[Ch92]{Ch92} H.~Chabanne, Permutation decoding of abelian codes,
\emph{IEEE Trans.\ Inform.\ Theory} \textbf{38} (1992), 1826--1829.

\bibitem[Du87]{Du87} A.~D\"ur, The automorphism groups of
Reed--Solomon codes, \emph{J.\ Combin.\ Theory Ser.~A} \textbf{44}
(1987), 69--82.

\bibitem[FR93]{FR93} G.-L.~Feng and T.~R.~N.~Rao, Decoding
algebraic-geometric codes up to the designed minimum distance,
\emph{IEEE Trans.\ Inform.\ Theory} \textbf{39} (1993), 37--45.

\bibitem[Go03]{Go03} N.~G\"ob, Computing the automorphism groups of
hyperelliptic function fields, arXiv:math/0305284.

\bibitem[Go81]{Go81} V.~D.~Goppa, Codes on algebraic curves,
\emph{Soviet Math.\ Dokl.} \textbf{24} (1981), 170--172.

\bibitem[Gor82]{Gor82} D.~M.~Gordon, Minimal permutation sets for
decoding the binary Golay codes, \emph{IEEE Trans.\ Inform.\ Theory}
\textbf{28} (1982), 541--543.

\bibitem[GS99]{GS99} V.~Guruswami and M.~Sudan, Improved decoding of
Reed--Solomon and algebraic-geometry codes, \emph{IEEE Trans.\ Inform.\
Theory} \textbf{45} (1999), 1757--1767.

\bibitem[GUAVA]{GUAVA} GUAVA, a GAP package for computing with
error-correcting codes, \url{https://www.gap-system.org/Packages/guava.html}.

\bibitem[Gu04]{Gu04} R.~Guralnick, private communication, 2004.

\bibitem[Ha77]{Ha77} R.~Hartshorne, \emph{Algebraic Geometry}, Springer, 1977.

\bibitem[HvLP98]{HvLP98} T.~H{\o}holdt, J.~H.~van Lint and R.~Pellikaan,
Algebraic geometry codes, in \emph{Handbook of Coding Theory}, vol.~1,
V.~S.~Pless and W.~C.~Huffman (eds.), Elsevier, 1998, 871--961.

\bibitem[HP03]{HP03} W.~C.~Huffman and V.~Pless, \emph{Fundamentals of
Error-Correcting Codes}, Cambridge Univ.\ Press, 2003.

\bibitem[Hu98]{Hu98} W.~C.~Huffman, Codes and groups, in \emph{Handbook
of Coding Theory}, vol.~2, Elsevier, 1998, 1345--1440.

\bibitem[Jo05]{Jo05} D.~Joyner, Conjectural permutation decoding of some
AG codes, \emph{ACM SIGSAM Bull.} \textbf{39}(1) (2005), 26--32,
DOI 10.1145/1080368.1080374.

\bibitem[Jo26]{Jo26} D.~Joyner et al., \texttt{perm-decoding}, folder of
the GitHub repository
\url{https://github.com/wdjoyner/qrcodes-hyperelliptic}, 2026.

\bibitem[JK06]{JK06} D.~Joyner and A.~Ksir, Automorphism groups of some
AG codes, \emph{IEEE Trans.\ Inform.\ Theory} \textbf{52} (2006),
3325--3329; arXiv:math/0412459.

\bibitem[JK03]{JK03} D.~Joyner and A.~Ksir, Representations of finite
groups on Riemann--Roch spaces, II, arXiv:math/0312383.

\bibitem[JKT08]{JKT08} D.~Joyner, A.~Ksir and W.~Traves, Automorphism groups
of generalized Reed--Solomon codes, arXiv:0801.4007.

\bibitem[JT02]{JT02} D.~Joyner and W.~Traves, Representations of finite
groups on Riemann--Roch spaces, arXiv:math/0210408.

\bibitem[JLJHH89]{JLJHH89} J.~Justesen, K.~J.~Larsen, H.~E.~Jensen,
A.~Havemose and T.~H{\o}holdt, Construction and decoding of a class of
algebraic geometry codes, \emph{IEEE Trans.\ Inform.\ Theory} \textbf{35}
(1989), 811--821.

\bibitem[KMM05]{KMM05} J.~D.~Key, T.~P.~McDonough and V.~C.~Mavron,
Partial permutation decoding for codes from finite planes,
\emph{European J.\ Combin.} \textbf{26} (2005), 665--682.

\bibitem[KMM06]{KMM06} J.~D.~Key, T.~P.~McDonough and V.~C.~Mavron,
Information sets and partial permutation decoding for codes from finite
geometries, \emph{Finite Fields Appl.} \textbf{12} (2006), 232--247.

\bibitem[KMM08]{KMM08} J.~D.~Key, T.~P.~McDonough and V.~C.~Mavron,
Partial permutation decoding for codes from affine geometry designs,
\emph{J.\ Geom.} \textbf{88} (2008), 101--109.

\bibitem[KMM10]{KMM10} J.~D.~Key, T.~P.~McDonough and V.~C.~Mavron,
Reed--Muller codes and permutation decoding, \emph{Discrete Math.}
\textbf{310} (2010), 3114--3119.

\bibitem[KMM17]{KMM17} J.~D.~Key, T.~P.~McDonough and V.~C.~Mavron,
Improved partial permutation decoding for Reed--Muller codes,
\emph{Discrete Math.} \textbf{340} (2017), 722--728.

\bibitem[KMR04]{KMR04} J.~D.~Key, J.~Moori and B.~G.~Rodrigues,
Permutation decoding for the binary codes from triangular graphs,
\emph{European J.\ Combin.} \textbf{25} (2004), 113--123.

\bibitem[KS07]{KS07} J.~D.~Key and P.~Seneviratne, Binary codes from
rectangular lattice graphs and permutation decoding, \emph{European J.\
Combin.} \textbf{28} (2007), 121--126.

\bibitem[KS08]{KS08} J.~D.~Key and P.~Seneviratne, Permutation decoding
for binary codes from lattice graphs, \emph{Discrete Math.} \textbf{308}
(2008), 2862--2867.

\bibitem[KW90]{KW90} R.~Kn\"orr and W.~Willems, The automorphism groups of
generalized Reed--Muller codes, \emph{Ast\'erisque} \textbf{181--182}
(1990), 195--207.

\bibitem[KV05]{KV05} H.-J.~Kroll and R.~Vincenti, PD-sets for the codes
related to some classical varieties, \emph{Discrete Math.} \textbf{301}
(2005), 89--105.

\bibitem[Le82]{Le82} J.~S.~Leon, Computing automorphism groups of
error-correcting codes, \emph{IEEE Trans.\ Inform.\ Theory} \textbf{28}
(1982), 496--511.

\bibitem[LLMS26]{LLMS26} M.~Lichtenwalner, H.~H.~L\'opez,
G.~L.~Matthews and P.~Seneviratne, Permutation decoding of algebraic
geometry codes from Hermitian and norm-trace curves, arXiv:2602.20455
(2026).

\bibitem[LFK10]{LFK10} F.~Lim, M.~Fossorier and A.~Kav\v{c}i\'c, Code
automorphisms and permutation decoding of certain Reed--Solomon binary
images, \emph{IEEE Trans.\ Inform.\ Theory} \textbf{56} (2010),
5253--5273.

\bibitem[Mac64]{Mac64} F.~J.~MacWilliams, Permutation decoding of
systematic codes, \emph{Bell System Tech.\ J.} \textbf{43} (1964),
485--505.

\bibitem[MS77]{MS77} F.~J.~MacWilliams and N.~J.~A.~Sloane, \emph{The
Theory of Error-Correcting Codes}, North-Holland, 1977.

\bibitem[Pr62]{Pr62} E.~Prange, The use of information sets in decoding
cyclic codes, \emph{IRE Trans.\ Inform.\ Theory} \textbf{8} (1962),
no.~5, 5--9.

\bibitem[Sc90]{Sc90} C.~Schoen, On certain modular representations in the
cohomology of algebraic curves, \emph{J.\ Algebra} \textbf{135} (1990),
1--18.

\bibitem[Sen09]{Sen09} P.~Seneviratne, Partial permutation decoding for
the first-order Reed--Muller codes, \emph{Discrete Math.} \textbf{309}
(2009), 1967--1970.

\bibitem[SF72]{SF72} S.~Shiva and K.~Fung, Permutation decoding of
certain triple-error-correcting binary codes, \emph{IEEE Trans.\ Inform.\
Theory} \textbf{18} (1972), 444--446.

\bibitem[SFT70]{SFT70} S.~Shiva, K.~Fung and H.~Tan, On permutation
decoding of binary cyclic double-error-correcting codes of certain
lengths, \emph{IEEE Trans.\ Inform.\ Theory} \textbf{16} (1970),
641--643.

\bibitem[ShW07]{ShW07} T.~Shaska and Q.~Wang, On the automorphism groups
of some AG-codes based on $C_{a,b}$ curves, \emph{Serdica J.\ Computing}
\textbf{1} (2007), 193--206.

\bibitem[SW99]{SW99} M.~A.~Shokrollahi and H.~Wasserman, List decoding
of algebraic-geometric codes, \emph{IEEE Trans.\ Inform.\ Theory}
\textbf{45} (1999), 432--437.

\bibitem[SV90]{SV90} A.~N.~Skorobogatov and S.~G.~Vl\u{a}du\c{t}, On
the decoding of algebraic-geometric codes, \emph{IEEE Trans.\ Inform.\
Theory} \textbf{36} (1990), 1051--1060.

\bibitem[St90]{St90} H.~Stichtenoth, On automorphisms of geometric Goppa
codes, \emph{J.\ Algebra} \textbf{130} (1990), 113--121.

\bibitem[St93]{St93} H.~Stichtenoth, \emph{Algebraic Function Fields and
Codes}, Springer, 1993 (2nd ed.\ 2009).

\bibitem[Su97]{Su97} M.~Sudan, Decoding of Reed--Solomon codes beyond
the error-correction bound, \emph{J.\ Complexity} \textbf{13} (1997),
180--193.

\bibitem[TV91]{TV91} M.~A.~Tsfasman and S.~G.~Vl\u{a}du\c{t},
\emph{Algebraic-Geometric Codes}, Kluwer, Dordrecht, 1991.

\bibitem[TVZ82]{TVZ82} M.~A.~Tsfasman, S.~G.~Vl\u{a}du\c{t} and
T.~Zink, Modular curves, Shimura curves, and Goppa codes, better than
Varshamov--Gilbert bound, \emph{Math.\ Nachr.} \textbf{109} (1982),
21--28.

\bibitem[We98]{We98} S.~Wesemeyer, On the automorphism group of various Goppa
codes, \emph{IEEE Trans.\ Inform.\ Theory} \textbf{44} (1998), 630--643.

\bibitem[Wo83]{Wo83} J.~Wolfmann, A permutation decoding of the
$(24,12,8)$ Golay code, \emph{IEEE Trans.\ Inform.\ Theory} \textbf{29}
(1983), 748--750.

\end{thebibliography}

\end{document}
